%% Chapter8_Conclusion.tex
\chapter{Conclusion and Future Scope}
\label{chap:conc}

\section{Summary}
This project addressed the problem that low-cost density-based signal controllers are rarely evaluated against strong baselines, near and above capacity, or for signal safety. The following was accomplished:
\begin{itemize}
  \item Designed and developed a complete simulation-first traffic management system: a microscopic simulator, a density-aware adaptive controller with a discharge hold, an emergency-preemption layer, a Haar-cascade and ByteTrack-style vision pipeline, a FastAPI backend with SQLite persistence, a React dashboard, a hardware bridge, Docker packaging and CI with 119 tests (Chapters~\ref{chap:meth} and~\ref{chap:impl}).
  \item Achieved a reduction of mean stopped delay below capacity of 35.3--36.1\% relative to a 20~s fixed plan, and with the discharge hold restored and exceeded the fixed plan's capacity (+9.8\%), as shown in the results chapter (Chapter~\ref{chap:results}).
  \item Demonstrated alignment with SDG 11, with secondary links to SDG 3, 9 and 13 (Chapter~\ref{chap:societal}).
\end{itemize}

\section{Conclusions}
Based on the experimental results and analysis, the following conclusions are drawn:
\begin{enumerate}
  \item Conclusion 1: Below capacity, the density-based controller reduces delay by 35.3--36.1\% relative to a 20~s fixed plan, but a Webster-optimized plan, actuated control and max-pressure control do as well; most of the gain comes from replacing an untuned plan, not from the density score.
  \item Conclusion 2: Near capacity, density-only switching ended greens too early and lost 18.9\% of capacity. A discharge-hold rule restored it (+9.8\% relative to fixed timing), lowered delay from 0.6 to 1.5 times capacity (not significantly at 1.0$\times$) and kept queues and spillback below those of fixed timing in a 60-minute peak.
  \item Conclusion 3: The queued/moving distinction, not the class weights, accounts for the benefit of the density score; data-driven calibration gave only modest, condition-dependent gains.
  \item Conclusion 4: Signal safety held in every run (no conflicting green or shortened change interval in 3,524 runs). The findings hold with kinematics calibrated to field saturation headways once the hold's zone is re-tuned (20.0~m).
  \item Conclusion 5: The present Haar-cascade detector (4.1\% precision, 8.2\% recall on synthetic frames) is too weak to drive control: with it, every vehicle-reading controller would be worse than fixed timing near capacity.
\end{enumerate}

\section{Original Contributions}
\label{sec:contrib}
The original contributions of this project are:
\begin{enumerate}
  \item A class-weighted, queue-aware density score with a bounded green-extension controller that never shortens yellow or all-red intervals and guarantees a maximum red time, a discharge-hold extension that removes the capacity loss of density-only switching, and an emergency-preemption layer that wraps any phase-based controller and keeps full change intervals, with a per-step signal audit.
  \item An empirical evaluation demonstrating, on paired seeded runs, how the controller compares with four re-implemented strategies below, at and above capacity, in a peak period, under calibrated kinematics and under imperfect sensing, including an ablation and a held-out calibration of the score.
  \item A publicly available codebase at \url{https://github.com/shivam-G05/Smart-Traffic-Management}, in which every table can be regenerated by a script, and a revised research paper reporting the results.
\end{enumerate}

\section{Limitations}
Results are from simulation only: the model has no turning movements, lane changes, pedestrians or driver variability, and arrivals are not Poisson. The main kinematic model has no driver reaction time (a calibrated variant is also reported), the baselines are re-implementations with ideal sensing, the vision pipeline was tested only on synthetic frames and is too weak to drive control, and the hardware bridge has not been built and tested on a physical board. The limitations were managed by calibrating the kinematics against the Highway Capacity Manual, by repeating the comparisons under the calibrated model, by measuring the effect of imperfect sensing, and by labelling every estimate as an estimate.

\section{Future Scope}
\label{sec:future}
Future work can extend this project in the following directions:
\begin{itemize}
  \item \textbf{Short-term:} improve resumption after preemption (return to the higher-scoring direction); add Poisson arrivals and turning movements; train a class-aware detector that suppresses static false positives near the stop line.
  \item \textbf{Medium-term:} close the loop between detection and control with that detector on real footage; test the controller in an established simulator such as SUMO~\cite{ref22}; build and verify the Raspberry Pi/ESP32 prototype and its timing on actual hardware; calibrate the class weights against local passenger-car-unit factors and field data.
  \item \textbf{Long-term:} pilot at a real junction with independent safety validation; extend to coordinated multi-junction corridors; study reinforcement-learning controllers with the same safety audit; broader deployment with municipal authorities.
\end{itemize}
